\documentclass[11pt]{article}

% ---- Packages ----
\usepackage[margin=1in]{geometry}
\usepackage{amsmath}
\usepackage{amssymb}
\usepackage{booktabs}
\usepackage{enumitem}
\usepackage{titlesec}
\usepackage[hidelinks]{hyperref}
\usepackage{xcolor}

% ---- Light formatting ----
\definecolor{moss}{HTML}{758C58}
\titleformat{\section}{\large\bfseries\color{moss}}{\thesection}{0.6em}{}
\titleformat{\subsection}{\normalsize\bfseries}{\thesubsection}{0.6em}{}
\setlist{itemsep=2pt, topsep=4pt}
\setlength{\parskip}{0.4em}
\setlength{\parindent}{0pt}

% ---- Title block ----
\title{\vspace{-1.5em}\textbf{Claude Prices Probability, Not Payoffs}\\[0.3em]
\large Phase 1 Results: Structural Estimation of Risk Preferences in an LLM}
\author{J. Blakeslee \\ \small Pardee RAND Graduate School}
\date{\small AI Sprint Practicum (Course 1012) \\ \small July 2026}

\begin{document}
\maketitle

\hrule
\vspace{0.5em}

\section{Introduction}

This project treats large language models as economic agents and recovers their preference
primitives from their choices. The full study elicits binary lottery choices from LLMs, estimates
structural discrete-choice models on the resulting data, and then carries the identical instruments
to matched base and instruction-tuned models to test whether post-training instills behavioral
biases. Phase 1 is the deliberate warm-up: fit a constant-relative-risk-aversion (CRRA) expected
utility model to 1{,}200 controlled choices from a single frontier model, and in doing so debug the
elicitation pipeline before the stakes rise. The purpose of a warm-up is to be easy. This one was
not, and the way it failed is the finding.

The model under study is \texttt{claude-haiku-4-5-20251001}, a pinned snapshot, queried at
temperature 1.0 through forced tool-choice with a 10\% plain-text validation arm. We asked it 1{,}200
either/or money questions of the form ``a sure \$86 or a 43\% chance of \$210,'' each one
procedurally generated to be contamination-safe, fully counterbalanced for option order, and drawn
under one of five paraphrase templates in a fresh context. A textbook CRRA fit returned a plausible
risk-aversion number and hid a specification failure underneath it. The diagnostics that the design
was built to run caught the failure. Once corrected, the model reveals a preference structure that is
coherent, human-implausible in a specific and quantifiable way, and directly relevant to how agentic
systems will behave when they allocate real resources.

\section{The Result in Plain Language}

We asked an AI model twelve hundred simple money questions. Every question had the same shape: a
guaranteed amount of money, or a gamble. For example, ``Would you rather take a sure \$86, or a 43\%
chance at \$210 and otherwise nothing?'' We varied the guaranteed amount, the size of the prize, and
the odds of winning across forty distinct gambles, asked each one many times with different wording
and different ordering, and recorded which side the model picked. Nothing was left to chance in the
design: the numbers were odd and non-round so the model could not be reciting a memorized textbook,
and every gamble appeared in both possible orders so we could tell whether mere position was steering
the answer.

A person deciding these questions weighs two things: how much money is on the table, and how likely
they are to win it. A long shot at a huge prize can be worth taking; a near-certain shot at a small
prize can be worth passing on. The model does something stranger. It looks almost entirely at the
odds of winning and barely at the size of the prize.

Here is the behavior in one picture. Offer the model a gamble it will probably lose, and it refuses,
no matter how large the prize. Raise the prize from \$200 to \$2{,}000 to \$20{,}000 while keeping
the win probability low, and the refusal does not budge. Now offer it a gamble it will probably win,
even one whose prize is only mediocre, and it takes it. When the chance of winning was below roughly
30\%, the model gambled about 8\% of the time. When the chance of winning was above roughly 60\%, it
gambled about 86\% of the time. The average money-value of the gambles was nearly identical across
those two groups. The model was not chasing value. It was chasing the likely winner.

This is close to the opposite of a well-known human quirk. People buy lottery tickets: they
overvalue tiny chances at enormous prizes. This model would never buy a lottery ticket at any price,
and it would happily take a coin-flip for pocket change. Neither is ``rational'' in the economist's
sense, and the model's version is arguably the more consequential one for anything asked to make real
decisions.

We found one more thing, and it is almost embarrassing in how simple it is. Merely listing the
gamble second, rather than first, made the model much more likely to choose it. Same gamble, same
numbers, same odds; only the position on the page changed, and the model's choices shifted by a large
and statistically overwhelming margin. Whoever controls the order in which options are presented has a
lever on what this model decides.

Why care? Because these models are being handed real decisions: which supplier to pick, which trade
to place, which case to prioritize, what to offer in a negotiation. A decision-maker that ignores how
much is at stake and fixates on the odds will systematically walk past high-value long shots (the
research bet, the insurance-like tradeoff) and grab safe mediocrity. That is a predictable, exploitable
distortion, and Phase 1 measures it.

\section{Experimental Design}

Every choice in the study is an independent API call with no memory of prior trials. The instrument
is a forced binary choice between a sure amount and a two-outcome gamble. The gamble pays a high
amount with probability $p$ and zero otherwise. The design holds to the measurement commitments the
proposal set out, treating the known validity threats in this literature as first-order.

\begin{itemize}
  \item \textbf{Contamination-safe gambles.} Every lottery is procedurally generated with non-round
  payoffs and jittered probabilities inside novel wording. There are no textbook Holt-Laury rows to
  recite.
  \item \textbf{Full counterbalancing of order.} Each gamble is presented in both orders, with the
  sure option and the gamble swapping first and second position, so position can be estimated rather
  than assumed away.
  \item \textbf{Template as a factor.} Each gamble is drawn under five semantically equivalent
  paraphrase templates, so between-template variance is observable.
  \item \textbf{Reproducibility.} Every trial logs the model snapshot, seed, temperature, and raw
  response. Elicitation is via forced tool-choice, with a 10\% plain-text arm to validate that the
  tool-choice channel is not itself distorting answers.
\end{itemize}

The design grid is 40 gambles (sure amount \$23--\$187 against a two-outcome gamble with high payoff
at probability $p \in [0.19, 0.83]$, else zero; EV ratios spanning 0.75--1.7) $\times$ 5 templates
$\times$ 2 orders $\times$ 3 repetitions, for 1{,}200 target choices. The final dataset contains
1{,}200 valid choices with zero unparseable responses. A mid-run API credit outage caused 265 trials
to fail; these were re-run cleanly, and because run order was shuffled, the outage is random with
respect to experimental conditions.

\section{Econometric Model and Estimation}

\subsection{Baseline: CRRA expected utility with a Fechner error (M1)}
The canonical structural exercise models the model as a CRRA expected-utility maximizer with a logit
(Fechner) choice error. Utility is
\[
  u(x) = \frac{x^{1-r}}{1-r},
\]
and the probability of choosing the gamble over the sure option is
\[
  \Pr(\text{gamble}) = \frac{1}{1 + \exp\!\big(-\,(EU_g - EU_s)/(\mu \cdot \text{scale})\big)},
\]
where $EU_g = p \cdot u(\text{hi})$, $EU_s = u(\text{sure})$, and $\text{scale} = |u(\text{max
payoff})|$ normalizes the noise term across stakes. Maximum likelihood on the 1{,}200 choices
returns $r = 0.19$ (se 0.048, 95\% CI $[0.10, 0.29]$) and $\mu = 0.26$ (se 0.040), with log-likelihood
$-809.69$. Read on its own, this is an unremarkable, mildly risk-averse agent.

\subsection{Diagnostic: probability drives choice, payoffs do not}
The sanity checks the design was built to run flag M1 as misspecified. The share of gambles chosen is
non-monotone across EV-ratio quartiles (0.43, 0.71, 0.45, 0.62); there is a large position gap of
0.12 ($p < 0.001$); and the spread across templates is 0.25. A descriptive logistic regression of the
gamble choice on the EV ratio, the win probability $p$, and a \texttt{safe\_first} indicator (the
gamble listed second) makes the pattern explicit:
\[
  \text{ev\_ratio: } +1.71\ (z=5.0), \quad p: +11.87\ (z=22.8), \quad \text{safe\_first: } +0.97\ (z=6.3).
\]
Holding the mean EV ratio roughly constant (between 1.08 and 1.19) across win-probability bins, uptake
climbs from 8.3\% to 85.9\% as $p$ rises (Table~\ref{tab:bins}). Choice tracks the probability of
winning; payoff magnitudes barely register. CRRA cannot generate this: it integrates probability and
payoff into a single expected utility, so it cannot make choice nearly independent of payoff while
sharply dependent on $p$.

\subsection{Extended model: Prelec weighting plus a position parameter (M2)}
The diagnostic points to a distorted probability weighting function. M2 adds a Prelec (1998) weighting
function and a position nuisance parameter $\delta$ that enters the latent index when the gamble is
listed second:
\[
  w(p) = \exp\!\big(-(-\ln p)^{\gamma}\big), \qquad
  z = \frac{w(p)\,u(\text{hi}) - u(\text{sure})}{\mu \cdot \text{scale}} + \delta \cdot \mathbf{1}[\text{gamble second}],
\]
with $\Pr(\text{gamble}) = 1/(1 + e^{-z})$. At $\gamma = 1$ and $\delta = 0$ this nests M1 exactly, so
the two are comparable by a likelihood-ratio test. Estimation is by maximum likelihood with
multi-start Nelder-Mead and standard errors from the numerical Hessian.

\section{Results}

Table~\ref{tab:params} reports the M2 estimates. Table~\ref{tab:compare} compares the baseline, the
extended model, and a bare-bones probability-only logit (M3). Table~\ref{tab:bins} shows fit against
the observed choice rate by win-probability bin.

\begin{table}[h]
\centering
\caption{M2 parameter estimates ($n = 1{,}200$).}
\label{tab:params}
\begin{tabular}{@{}lll@{}}
\toprule
Parameter & Estimate (se) & Interpretation \\
\midrule
$r$ (curvature)            & 0.592 (0.028) & Risk aversion; human-plausible range \\
$\gamma$ (Prelec)          & 2.693 (0.166) & Extreme S-shape; long shots underweighted \\
$\mu$ (Fechner noise)      & 0.104 (0.008) & Choice noise, less than half of M1's \\
$\delta$ (position)        & 1.257 (0.160) & Shift toward the gamble when it is listed second \\
\bottomrule
\end{tabular}
\end{table}

\begin{table}[h]
\centering
\caption{Model comparison. Lower AIC/BIC is better.}
\label{tab:compare}
\begin{tabular}{@{}lccrrr@{}}
\toprule
Model & Structure & $k$ & logLik & AIC & BIC \\
\midrule
M1 & CRRA + Fechner            & 2 & $-809.69$ & 1623.4 & 1633.6 \\
M2 & CRRA + Prelec + position  & 4 & $-501.24$ & 1010.5 & 1030.8 \\
M3 & Probability-only logit    & 2 & $-501.40$ & 1010.8 & 1021.0 \\
\bottomrule
\end{tabular}
\end{table}

The likelihood-ratio test of M1 against M2 is decisive: $\chi^2(2) = 616.9$, $p = 1.1 \times
10^{-134}$. Adding probability weighting and a position term transforms the fit.

\begin{table}[h]
\centering
\caption{Choice rate by win-probability bin: observed vs.\ fitted. EV ratio is roughly constant across
bins, so a value-integrating model should be near-flat.}
\label{tab:bins}
\begin{tabular}{@{}lccc@{}}
\toprule
Win probability $p$ & Observed & M1 fitted & M2 fitted \\
\midrule
$(0.15, 0.30]$ & 0.083 & 0.427 & 0.028 \\
$(0.30, 0.45]$ & 0.138 & 0.474 & 0.161 \\
$(0.45, 0.60]$ & 0.607 & 0.493 & 0.654 \\
$(0.60, 0.85]$ & 0.859 & 0.487 & 0.782 \\
\bottomrule
\end{tabular}
\end{table}

M1 predicts a nearly flat 0.43--0.49 across bins, missing the observed climb entirely. M2 tracks the
step from near-zero uptake at low $p$ to high uptake at high $p$.

\section{Interpretation}

\textbf{The weighting function is extreme, and points the wrong way.} A Prelec $\gamma = 2.69$ is a
strong S-shape. It severely underweights low and moderate probabilities and only credits near-certain
odds. Humans show the opposite: an inverse-S with $\gamma$ around 0.65--0.75 (Prelec 1998; Wu and
Gonzalez 1996), overweighting long shots, which is why people buy lottery tickets. Haiku radically
underweights anything short of a likely win. It will not take a long shot at nearly any price.

\textbf{Once weighting is controlled, the curvature is ordinary.} The M2 estimate $r = 0.59$ sits
squarely in the standard human CRRA range. The baseline $r = 0.19$ was misspecification bias: forcing
a probability-distortion pattern through a pure-curvature model pulled the curvature estimate toward
risk neutrality. The apparent risk attitude in M1 was an artifact of the wrong model.

\textbf{Position moves choice by a large factor.} A position parameter $\delta = 1.26$ implies an odds
shift of $e^{1.26} \approx 3.5\times$ toward the gamble merely from listing it second. The direction is
notable: this is a second-option bias, opposite the first-option ``A-bias'' that much of the LLM
literature reports.

\textbf{Noise was mostly misspecification.} The Fechner noise term halved, from $\mu = 0.26$ in M1 to
$\mu = 0.10$ in M2. What looked like random choice under the wrong model was largely structure the
wrong model could not represent.

\textbf{The honest caveat.} The probability-only heuristic M3 essentially ties M2 on AIC (1010.8 vs.\
1010.5) and beats it on BIC (1021.0 vs.\ 1030.8). Win probability alone explains almost everything;
payoff magnitudes add little. This near-tie is itself the finding. The model behaves closer to a
``pick the likely winner'' heuristic than to genuine expected-utility integration. The structural
model earns its place by interpretability: it nests the rational benchmark, yields parameters
comparable to the human behavioral literature, and locates exactly where and how the model departs
from that benchmark, which a bare logit cannot do.

\section{Policy Relevance}

\textbf{Allocative agents inherit this distortion.} LLMs are being deployed to make allocative choices:
procurement, trading, triage, negotiation. An agent that thresholds on probability and ignores payoff
magnitude will systematically decline high-value long shots (R\&D bets, insurance-like tradeoffs) and
accept likely-but-mediocre options. Phase 1 quantifies that distortion for one snapshot under one
protocol, turning a plausible worry into a measured parameter.

\textbf{Interface order is a control lever.} A $3.5\times$ position effect means whoever controls the
order in which options are presented steers the agent's choice. Interface design becomes a lever over
AI decisions. That matters for anyone building or regulating agentic systems, and it matters more in
adversarial settings, where an opponent who arranges the menu can tilt the outcome.

\textbf{Structural estimation earns its keep.} A descriptive bias demonstration cannot separate ``the
model is risk-averse'' from ``the model ignores payoffs.'' The M1 fit would have reported the former;
the truth is the latter. Structural estimation with a correctly specified model can tell them apart,
and the misspecification was only detectable because the design carried sanity checks and
counterbalancing. This is the argument for measurement-grade evaluations rather than one-off bias
anecdotes.

\textbf{Scope.} These parameters are properties of this pinned snapshot under this protocol:
unincentivized, forced-choice, one model. They are not claimed as universal traits of LLMs. Phase 2
(gain/loss framing, loss aversion, an endogenous reference point) and Phase 3 (base versus instruct
attribution) test whether the pattern is stable and where it comes from.

\section{Limitations and Next Phase}

Phase 1 studies a single model, at a single temperature, with unincentivized forced choice, so the
estimates are conditional on that protocol. The near-tie between the structural model and the
probability-only heuristic means the payoff-magnitude channel is weak in this design; a stakes range
wider than \$23--\$187 might surface more curvature, and Phase 2 will widen it. The position effect is
estimated as a single nuisance term rather than modeled by template, though the counterbalanced design
ensures it does not bias the preference parameters. Temperature 1.0 sampling is treated as a choice
error process here; whether that mapping holds is a separate question the logprob elicitation planned
for Phase 3 will address directly.

Phase 2 adds the gain/loss frame manipulation and estimates loss aversion and an endogenous reference
point. Phase 3 runs the identical instruments on matched base and instruction-tuned open-weight models
to attribute the estimated parameters to post-training. Phase 1 has done its job: the pipeline is
validated end to end, and it has already returned a sharp, quantified, and policy-relevant result.

\vspace{1em}
\hrule
\vspace{0.5em}

\begin{thebibliography}{99}
\small

\bibitem{holt2002}
Holt, C. A., and Laury, S. K. (2002).
Risk Aversion and Incentive Effects. \textit{American Economic Review}, 92(5), 1644--1655.

\bibitem{kahneman1979}
Kahneman, D., and Tversky, A. (1979).
Prospect Theory: An Analysis of Decision under Risk. \textit{Econometrica}, 47(2), 263--291.

\bibitem{koszegi2006}
K\H{o}szegi, B., and Rabin, M. (2006).
A Model of Reference-Dependent Preferences. \textit{Quarterly Journal of Economics}, 121(4), 1133--1165.

\bibitem{liu2025}
Liu, J., Yang, Y., and Tam, K. Y. (2025).
Evaluating and Aligning Human Economic Risk Preferences in LLMs.
\textit{Proceedings of EMNLP 2025}; arXiv:2503.06646.

\bibitem{cpt2025}
CPT Estimation on LLMs (2025). Cumulative Prospect Theory Parameters of Instruction-Tuned Language
Models. arXiv:2508.08992.

\bibitem{prelec1998}
Prelec, D. (1998).
The Probability Weighting Function. \textit{Econometrica}, 66(3), 497--527.

\bibitem{train2009}
Train, K. E. (2009).
\textit{Discrete Choice Methods with Simulation} (2nd ed.). Cambridge University Press.

\bibitem{wu1996}
Wu, G., and Gonzalez, R. (1996).
Curvature of the Probability Weighting Function. \textit{Management Science}, 42(12), 1676--1690.

\end{thebibliography}

\end{document}
